\documentclass[10pt]{article}
\usepackage[margin=20mm]{geometry}
\usepackage{amsmath,amssymb,mathtools,fontspec,unicode-math,xurl,hyperref}
\IfFontExistsTF{DejaVu Serif}{\setmainfont{DejaVu Serif}}{\setmainfont{Latin Modern Roman}}
\setmathfont{latinmodern-math.otf}
\setlength{\parindent}{0pt}\setlength{\parskip}{6pt}\sloppy
\title{Well-posedness of local potentials
in a polytopal exterior calculus}
\author{AI-assisted technical review}
\date{}
\begin{document}
\maketitle
Abstract: The VEM potential reconstruction in the reviewed manuscript is not well defined: its prescribed wedge pairing is degenerate, and legitimate lowest-order data on a triangle make its defining equation inconsistent. An interval already gives nonuniqueness. This Major issue affects the potential and its polynomial-consistency theorem, whereas the separately defined VEM cochain complex and its cohomology argument do not use the defective reconstruction. Five local corrections concern the zero-degree Hodge Laplacian, the boundary scope of adjoint consistency, a polynomial degree, and two signs. The DDR algebraic proof chain receives positive checks. These conclusions are an initial review of an exact author-manuscript version; independent acceptance remains pending.\par
\section*{Source and scope}
The primary source is Francesco Bonaldi, Daniele Antonio Di Pietro, Jérôme Droniou and Kaibo Hu, An exterior calculus framework for polytopal methods, arXiv:2303.11093v3, 18 January 2025, 55 pages [main]. Printed and one-based PDF page numbers agree. All 55 text pages were read and all 55 source page images actually viewed. The publisher identifies DOI 10.4171/JEMS/1602, acceptance on 10 November 2024 and online publication on 17 March 2025 [publisher]. The acquired v3 carries a submitted-manuscript header; equivalence with the journal PDF has not been established.\par
The review checked the Koszul decompositions and traces, both discrete-space definitions, local and global operators, complex and commutation arguments, polynomial consistency, cohomology reductions, the conditional Hodge-Laplacian application, and the appendix’s conventions. The external analytic estimates cited as [34], [36], [44] and [47] in the paper were not independently re-proved in their entirety. In particular, the paper explicitly leaves a unified general-degree adjoint-consistency proof to future work (p. 21); this acknowledged limitation is not treated as a proved error.\par
\section*{The VEM potential is not well defined}
Problem 1 (Major).  Equation (4.9), p. 37, does not in general define a unique element of \(\mathcal P^-_{r+1}\Lambda^k(f)\), and for permitted data it need not have a solution. The inference from moments to equality in the proof of Theorem 31, p. 40, uses the same invalid nondegeneracy assertion.\par
The test forms \(\mathrm d\mu+\nu\) in (4.9) span \(\mathcal P^-_{r+1}\Lambda^{d-k}(f)\) by (2.17). That fact does not imply that \[(P,\rho)\longmapsto\int_f P\wedge\rho,\qquad
 \mathcal P^-_{r+1}\Lambda^k(f)\times
 \mathcal P^-_{r+1}\Lambda^{d-k}(f)\longrightarrow\mathbb R
                                                        \tag{1}\] is a perfect pairing. The Hodge star does not preserve trimmed spaces in the way this conclusion would require.\par
\subsection*{Nonuniqueness on an interval}
Take \(f=(-1,1)\), \(d=1\), \(k=r=0\), and \(x_f=0\). Give all endpoint and cell components in (4.1) the value zero. The derivative (4.4), the boundary potentials (4.8), and \(\widetilde D\) in (4.10) are then zero. The trial space in (4.9) is \[\mathcal P^-_1\Lambda^0(f)=\operatorname{span}\{1,x\},\] whereas \(\mathcal K^0_1(f)=\operatorname{span}\{x\}\) and \(\mathcal K^1_1(f)=0\). Thus (4.9) requires only \(\int_{-1}^{1}P(x)\,\mathrm dx=0\). Both \(P=0\) and \(P=x\) satisfy every test. No endpoint matching condition is imposed on this reconstructed potential: the separately reconstructed boundary values merely enter its right-hand side.\par
\subsection*{Nonexistence on a regular triangle}
A stronger example avoids relying only on a dimension mismatch. Let \[T=\{x>0,\ y>0,\ x+y<1\},\quad
 d=2,\quad k=1,\quad r=0,\quad x_T=(1/3,1/3),\] and set \(X=x-1/3\), \(Y=y-1/3\) and \[\alpha=-Y\,\mathrm dx+X\,\mathrm dy=\kappa_T(\mathrm dx\wedge\mathrm dy)
 \in\mathcal K^1_1(T).\] Set every edge component to zero, the cell component to \(\omega_T=\alpha\), and \(D_{\omega,T}=0\). These are legal independent components of the Cartesian space (4.1): on this \((k+1)\)-cell its factor is \(\mathcal K^1_1(T)\times\mathcal K^0_0(T)\), with the second space zero. Equation (4.4) gives \(\mathrm d_{0,T}^1\underline\omega_T=0\), and the edge potentials are zero.\par
Here \[\mathcal P^-_1\Lambda^1(T)=\operatorname{span}\{\mathrm dx,\mathrm dy,\alpha\}.\] For every allowed trial \(P\), the centroid identities give \[\int_T P\wedge\alpha=0.\] Indeed, \(\int_T X=\int_T Y=0\) and \(\alpha\wedge\alpha=0\). Testing (4.9) with \((\mu,\nu)=(0,\alpha)\) instead requires \[0=\int_T\star^{-1}\alpha\wedge\alpha
   =\int_T(X^2+Y^2)\,\mathrm dx\,\mathrm dy
   =\frac1{18}.                                         \tag{2}\] The sign here follows from the paper’s own Hodge convention: \(\star^{-1}\alpha=X\,\mathrm dx+Y\,\mathrm dy\). Each centered second moment is \(1/36\), obtained from \(\int_Tx=1/6\), \(\int_Tx^2=1/12\), and \(|T|=1/2\). Thus the defining linear system is inconsistent even on a convex shape-regular cell.\par
\subsection*{A three-dimensional check and the repair boundary}
The same obstruction occurs for \(f=(-1,1)^3\), \(d=3\), \(k=2\), \(r=0\). Take all face components zero, \(\omega_f=x\,\mathrm dy-y\,\mathrm dx\) and \(D_{\omega,f}=0\). The proxy of any trial in \(\mathcal P^-_1\Lambda^2(f)\) has the form \(a+b(x,y,z)\). Its integral pairing with the proxy \((-y,x,0)\) vanishes, by oddness and pointwise radial orthogonality. Equation (4.9) at \((0,\omega_f)\) demands instead \(0=-16/3\). This is the squared-norm integral with the factor \((-1)^{k+1}=-1\).\par
This last datum is even the interpolant (4.3) of a smooth polynomial two-form. Put \(\phi=(1-x^2)(1-y^2)(1-z^2)\) and \[\beta=\frac{45}{8}\phi\,\star(x\,\mathrm dy-y\,\mathrm dx).\] Its face traces vanish. Symmetry makes its projected cell component a multiple of \(\omega_f\), and the multiplier is \[\frac{45}{8}
 \frac{\int_f\phi(x^2+y^2)}{\int_f(x^2+y^2)}
 =\frac{45}{8}\frac{128/135}{16/3}=1.\] The derivative component projects to \(\mathcal K^0_0=0\). The counterexample therefore does not depend on a possible objection that arbitrary product data lack a smooth antecedent.\par
Impact. The claimed map on the whole VEM space, the potential identity (4.12), and constructions using this map, including the higher-dimensional derivative (4.11), require substantive redesign. The last step of the proof of Theorem 31 only proves equality of moments in the degenerate pairing (1). For polynomial inputs, satisfying those moments alone cannot select the asserted unique polynomial.\par
Counterevidence. The global differential (4.6) uses (4.4) and the stored derivative components, not (4.9). Lemma 32, Propositions 33–34, and the cohomology argument in Theorem 30 survive this particular objection. The potentials in its extension (4.20) are DDR(0) potentials. No counterexample to the VEM cohomology theorem or to the DDR complex is asserted here.\par
Repair. A nondegenerate trial–test pair and a compatible polynomial-reproduction argument are needed. Simply replacing the target by \(\star^{-1}\mathcal P^-_{r+1}\Lambda^{d-k}\) would restore a Riesz pairing but does not preserve the theorem: in the interval example the new target is constants, so the interpolant of \(x\) could reconstruct only its mean rather than \(x\). No complete alternative VEM construction is certified in this review.\par
\section*{The Hodge-Laplacian application’s scope}
Problem 2 (Minor).  Section 3.2.1, p. 18, includes \(k=0\) while inferring that all harmonic spaces vanish from topological triviality. Consequently the unrestricted uniqueness claim and the inf-sup inequality (3.21), p. 19, fail at that endpoint.\par
On a nonempty connected contractible domain the unreduced de Rham complex used in (1.2) has \(H^0\cong\mathbb R\). Let \(c_h=I_r^0 1\). Polynomial consistency and commutation give \[c_h\ne0,\qquad \mathrm d_r^0c_h=0.\] Even with the usual convention \(X_r^{-1}=0\), \[\mathcal A_h((0,c_h),(0,v_h))=0
 \quad\hbox{for every }v_h,\qquad
 \lVert(0,c_h)\rVert_{\mathrm{graph}}>0.\] This contradicts (3.21). For zero forcing there are multiple constant solutions; for \(g=1\), testing with \(c_h\) gives zero on the left and \((c_h,c_h)_{0,h}>0\) on the right.\par
Restricting the displayed application to \(1\le k\le n\) on a contractible domain repairs this endpoint. Alternatively, impose zero mean and compatible forcing at \(k=0\), or retain the harmonic unknown and the usual orthogonality condition. Remark 21 already describes the relevant kind of harmonic correction. This does not refute (3.20): a constant input has zero differential and admits the chosen preimage zero. Nor does it refute Theorem 14, whose cohomology statement correctly retains constants.\par
Problem 3 (Minor).  The universal smooth-form scope of the adjoint-consistency estimate (3.24), p. 20, omits the boundary condition needed to remove the continuous boundary pairing.\par
Take \(\Omega=(0,1)^3\), \(\ell=0\), \(r\ge1\), a regular refining mesh, \(\omega=\mathrm dx_1\), and \(\mu_h=I_r^0x_1\). Then \(\delta\omega=0\) and \(\mathrm d_r^0\mu_h=I_r^1\mathrm dx_1\). Polynomial consistency, (3.16), and vanishing stabilization on polynomial interpolants give exactly \[\widetilde E_{r,h}^0(\mathrm dx_1;I_r^0x_1)=-1,\qquad
 \lVert I_r^0x_1\rVert_{\mathrm{graph}}^2=\frac13+1=\frac43. \tag{3}\] No constant independent of \(h\) can make (3) bounded by \(Ch^{r+1}\sqrt{4/3}\) as \(h\to0\). The forms are smooth, so additional piecewise regularity alone cannot exclude the example.\par
The statement should specify the Hilbert-adjoint domain, with the corresponding homogeneous normal trace (in smooth notation, \(\operatorname{tr}_{\partial\Omega}(\star\omega)=0\)), or explicitly retain the boundary term. The weak mixed equations can enforce the appropriate natural conditions for the actual solution \(u\) and \(\mathrm du\). Thus this counterexample targets the stated general estimate, not the conditional convergence argument after its boundary assumptions are restored. The paper’s acknowledged dependence on external Poincaré and adjoint estimates is separately retained as a coverage limitation.\par
\section*{Local formula corrections}
Problem 4 (Minor).  The two-dimensional radial proxy in (2.11), p. 7, has \(\mathcal P_{r+1}(F)\) where \(\mathcal P_{r-1}(F)\) is required.\par
By (2.5), \(\mathcal K_r^1(F)=\kappa_F\mathcal P_{r-1}\Lambda^2(F)\). After the specified rotation its vector proxy is therefore \[\mathcal R_r^c(F)=(x-x_F)\mathcal P_{r-1}(F).\] At \(r=0\), (2.6) makes the left-hand space zero, while the printed \((x-x_F)\mathcal P_1(F)\) is nonzero. Correcting the degree also corrects the explicit proxy interpretation in Example 3, Table 1, Example 13 and Table 2. The formal Koszul and trimmed-space definitions used in the algebraic DDR proofs have the correct degree; those proofs are not invalidated by this local proxy formula.\par
Problem 5 (Minor).  The second bullet following (A.6), p. 50, reverses the cross-product order for contraction of a two-form.\par
With the proxy convention on p. 49, take \(\omega=\mathrm dx_1\wedge\mathrm dx_2\), whose proxy is \(w=e_3\), and \(v=e_1\). The first-slot definition (A.6) gives \[\iota_{e_1}(\mathrm dx_1\wedge\mathrm dx_2)=\mathrm dx_2
 \longleftrightarrow e_2,\] whereas the printed \(v\times w=e_1\times e_3=-e_2\). The correct proxy is \(w\times v\). The formal definition of contraction, and hence the formal Koszul spaces, remains usable; changing the sign of a space generator does not change its span. Implementations translating the operator through this bullet need the correction.\par
Problem 6 (Minor).  The unnumbered integration-by-parts reformulation of (3.4), p. 47, has the wrong sign on its volume derivative.\par
Write \(a=\star^{-1}\omega_f\) and let \(P_{\partial f}\) denote the boundary potential. Stokes’ formula (2.1) and (3.4) give \[\int_f\mathrm d_r^k\underline\omega_f\wedge\mu
 =\int_f\mathrm da\wedge\mu+
   \int_{\partial f}(P_{\partial f}-\operatorname{tr}a)\wedge\operatorname{tr}\mu.   \tag{4}\] The printed first term is its negative. For a concrete check, take the interval \((0,1)\), \(k=0\), \(r=2\), and DDR data \(I_2^0x\). The cell component has \(a=x\), the endpoint potentials equal its traces, and the reconstructed derivative is \(\mathrm dx\). Testing with \(\mu=1\) makes (4) read \(1=1\); the printed formula reads \(1=-1\). Replacing the volume minus sign by a plus sign is a complete local repair. The original (2.1), (3.4) and the proved complex property use the correct sign and are not refuted by this explanatory rewrite.\par
\section*{Proof coverage, counterevidence and remaining limits}
The following positive checks delimit the adverse conclusions. For DDR, the full polynomial wedge pairing in (3.4) is nondegenerate; the Koszul decomposition (2.9) supplies the full test space in (3.5). The trace lemma accounts for the translation between cell centers. Boundary-of-boundary cancellation, with the orientation factors of (3.6), proves Lemma 22. Theorem 23 then establishes the complex property and the potential–derivative relation by induction on \(d-k\). Lemma 24 preserves the trimmed moments, Theorem 25 establishes commutation, and Theorem 15 uses these already proved relations in its polynomial-consistency induction. The \(r=0\) and \(d=k\) base cases were checked separately. These arguments do not invoke the defective VEM pairing.\par
For DDR cohomology, reduction \(R\) takes the lowest-dimensional component means. Constant tests in (3.56)–(3.58) show that it is a surjective cochain map. Its kernel is exactly the zero-mean complex, which Lemma 27 proves acyclic by constructing moments dimension by dimension. Consequently \[0\longrightarrow\ker R\longrightarrow\operatorname{DDR}(r)
   \longrightarrow\operatorname{DDR}(0)\longrightarrow0\] induces a cohomology isomorphism. The lowest-order complex is the cellular cochain complex after volume scaling. This independently checks the central reduction without erasing \(H^0\).\par
For VEM, Lemma 35 explicitly constructs a primitive in the kernel of the analogous mean reduction. Its lowest-dimensional compatibility uses zero integral, not a false identification of Koszul vanishing at a point with zero mean. Equations (4.19)–(4.23) give the reduction and extension through DDR(0), independently of (4.9). Theorem 30 therefore has a proof chain not contradicted by Problem 1.\par
Several smaller scope and notation points are retained without inflating the problem count. Equation (4.11), p. 37, defines the higher-dimensional derivative only for \(d\le n-1\), whereas (4.13) includes \(d=n\); the intended upper bound should be \(n\), subject to repairing the potential. The Hodge-space statement in Remark 5, p. 11, should distinguish \(\star^{-1}\mathcal P_r^-\Lambda^{d-k}\) from \(\mathcal P_r^-\Lambda^k\). For instance, in dimension three at \(r=1\) the latter has dimension six when \(k=1\), while \(\mathcal P_1^-\Lambda^2\) has dimension four. The malformed space membership following (3.39), p. 25, can be read correctly using the Hodge-transformed trimmed space; the moment argument then remains valid. Index slips in the reduction–extension prose were not promoted into false cohomology counterexamples.\par
The approximation arguments on pp. 27–28 were checked as reductions to the cited projector and stability estimates, rather than as new independent proofs of every mesh-uniform constant. No numerical convergence experiment or author implementation was run. No article-specific executable supplement was established by the bounded title/DOI, publisher, arXiv and repository searches. The Oxford repository lists an accepted manuscript; it was not silently substituted for the acquired v3. Absence of code is not a scientific defect. The current arXiv history still lists v3 as its latest version [main]; no journal-version correction is inferred from that fact.\par
\section*{Supporting checks}
The initial assessment is Major, because the stated VEM potential map fails to exist for valid input. This does not require the entire paper, either cohomology theorem, or every numerical application to be false. The other five numbered issues admit local corrections with the scope indicated above. A complete replacement reconstruction has not been proved here.\par
Three existing native helpers supplied bounded text-based checks. Their actual complete ranges were: the concentration-named helper, pp. 1–28; the waiting-time-named helper, pp. 1–13 and 15–38; and the mean-projection-named helper, pp. 1–14 and 35–55. Their additional partial reads and the fact that none viewed source images are recorded in provenance. The primary executor independently reconstructed every adopted example, read the complete source, viewed all source pages, wrote this report and performed its build and report-page inspection. Helper identities are task names, not evidence of a particular model or independent acceptance.\par
The primary PDF is 652856 bytes, with SHA256 \[\text{\small{7106d1b8fada237dbb2c4b4564dbb308ee278b5c3224b839bf2e1bb1b5aa36d3}}.\] Root-authored exact rational checks of (2), the cube integrals and (3) are retained with the actual source hashes. Compilation, dependency, cache and page-image identities and actual page inspections are bound in the accompanying records. No author code or shell escape was used. The controller owns the reported UI observation. Runtime self-description is recorded separately, with no assertion of independently verified backend equivalence. Scientific decision remains `pending`, human reviews remain empty, and independent acceptance is pending.\par
Conclusion.\par
Item | State\par
Supported findings | One Major VEM reconstruction failure and five Minor findings. DDR algebra and VEM cohomology remain separately supported.\par
Repair verification | Local degree/sign repairs verified; endpoint/boundary remedies bounded. No complete replacement VEM reconstruction or full analytic reproof is certified.\par
Author adoption | Not established for any of the six findings.\par
Remaining work | Construct a nondegenerate VEM potential and prove polynomial reproduction and the dependent reconstruction claims.\par
Scientific decision | Pending; human reviews empty; independent acceptance pending.\par
9 F. Bonaldi, D. A. Di Pietro, J. Droniou and K. Hu, An exterior calculus framework for polytopal methods, arXiv:2303.11093v3, 18 January 2025. https://arxiv.org/abs/2303.11093v3. EMS Press, article metadata for DOI 10.4171/JEMS/1602, https://ems.press/journals/jems/articles/14298659. Publisher PDF not collated with the reviewed author manuscript.\par
\end{document}
